\chapter{Reduced geometry and noncollapsing}
\label{ch:reduced-geometry}

This chapter isolates the analytic mechanism used later at high-curvature
points.  It starts with an ordinary Ricci flow on one fixed manifold, where
minimizing backward curves can be constructed by a direct method.  The same
objects are then rebuilt on a generalized spacetime.  The carrier is retained
through the gauge and rescaling constructions; a change of coordinates is not
silently treated as a new flow.

Throughout, $n$ is a positive integer, $M$ is a connected $n$-manifold with
the Hausdorff, regular and second-countable hypotheses, and $g(t)$ is a smooth
Ricci flow on an interval $J$.  Measures use the Borel sigma algebra.  In the
ordinary-flow sections $M$ may have arbitrary dimension; the compact
noncollapsing theorem will explicitly specialize to dimension three.  We use
\[
  \partial_tg=-2\operatorname{Ric},\qquad R=\operatorname{tr}_g\operatorname{Ric},
  \qquad |\mathrm{Rm}|=\text{the full curvature norm}.
\]
The base time is $T\in J$ and the backward time is $\tau\geq0$; the slice is
the metric $g(T-\tau)$.  A bound on $[T-\tau_{\max},T]$ always means a bound on
the full curvature tensor, not merely on $R$.

\section{Backward action and minimizing curves}

For a $C^1$ curve $\gamma:[\tau_1,\tau_2]\to M$ put
\[
 \mathcal L_{T}(\gamma)=\int_{\tau_1}^{\tau_2}
 \sqrt{\tau}\bigl(R(\gamma(\tau),T-\tau)
       +|\dot\gamma(\tau)|_{g(T-\tau)}^2\bigr)\,d\tau .
\]
An admissible backward path has time-window membership, continuity on the
closed interval, interior $C^1$ regularity and an integrable action density.  These
fields matter at $\tau=0$: they are not a convention that every continuous curve
is admissible.

\begin{definition}[Reduced length]
\label{def:reduced-length}
\lean{PoincareMT.reducedLength,PoincareMT.BackwardTimePath}
For $p,q\in M$ and $\tau>0$, define
\[
 \ell_T(p,q,\tau)=\frac1{2\sqrt\tau}
   \inf\{\mathcal L_T(\gamma):\gamma(0)=p,\ \gamma(\tau)=q\}.
\]
The infimum is over admissible backward paths on $[0,\tau]$.  A path is minimizing
when its action is no larger than that of every competitor with the same two
endpoints.
\end{definition}

\begin{theorem}[Existence, Euler equation and variation]
\label{thm:m08-l-geodesics}
\uses{def:reduced-length,prop:rf-evolution}
\lean{PoincareMT.lGeodesicExistenceAndVariation}
Let $T\in J$, $\tau_{\max}>0$, $[T-\tau_{\max},T]\subset J$, and assume a
complete full-curvature bound on that closed slab: each slice is complete
and one constant $C<\infty$ bounds $|\mathrm{Rm}|$ at every point and time
in the slab.  Then for every
$0\leq\tau_1<\tau_2\leq\tau_{\max}$ and every pair $p_1,p_2\in M$ there is a
minimizing backward path from $p_1$ at $	au_1$ to $p_2$ at $	au_2$.
Every minimizer satisfies the Euler--Lagrange equation on $(\tau_1,\tau_2)$,
and every Euler path has a square-root-time smooth representative at both
endpoints.  In particular the infimum defining $\ell_T$ is attained.
For a smooth square-root-time variation, the first derivative is the two
endpoint momentum terms plus the integral of the Euler residual.  For an
Euler path the second derivative is the endpoint acceleration term plus the
index form; when both endpoints are fixed and the path is minimizing, the
kernel of this form is exactly the space of endpoint-zero
$\mathcal L$-Jacobi fields.  A
positive-time Euler path extends uniquely to $\tau=0$, and the zero-start
Jacobi initial-value problem has a unique solution for each initial tangent
derivative.
\source{MT2007}\dcref{Definitions 6.1--6.2, p. 106}
\dcref{Lemmas 6.4, 6.8, 6.10, 6.12, pp. 107--110}
\dcref{Proposition 6.13, pp. 110--112}
\dcref{Proposition 6.33 and Claim 6.34, pp. 120--122}
\dcref{Lemmas 7.2--7.3, pp. 150--151}\cite[Definitions 6.1--6.2, Lemmas 6.4, 6.8, 6.10, 6.12, Proposition 6.13, Proposition 6.33, Claim 6.34, and Lemmas 7.2--7.3]{MT2007}
\end{theorem}
\begin{proof}
Write $s_i=\sqrt{\tau_i}$ and $\alpha(s)=\gamma(s^2)$.  Substitution removes
the degeneracy of the kinetic term:
\begin{equation}\label{eq:regularized-action}
 \mathcal L_T(\gamma)=\int_{s_1}^{s_2}
 \left(\tfrac12|\alpha'(s)|_{g(T-s^2)}^2
             +2s^2R(\alpha(s),T-s^2)\right)\,ds.
\end{equation}
The curvature bound gives a scalar bound and, by integrating
$|\partial_tg(v,v)|\leq C_n Cg(v,v)$, uniform equivalence with $g(T)$.
Consequently an upper action bound gives a uniform bound for
$\int|\alpha'|_{g(T)}^2$.  Cauchy--Schwarz gives
$d_{g(T)}(\alpha(a),\alpha(b))\leq C_1|a-b|^{1/2}$ and a common bounded
ball about the initial endpoint.  Completeness and Hopf--Rinow make this
ball compact.  A uniformly convergent subsequence therefore exists.
One first takes the limit in the absolutely continuous finite-energy class:
in finitely many charts the derivatives converge weakly in $L^2$, and the
positive quadratic kinetic integral is lower semicontinuous.  The scalar
term converges by uniform convergence on the compact spacetime image.
Smooth approximation of finite-energy chart curves preserves the endpoints
and action, so this enlarged minimization problem has the same infimum.

For clarity, the first variation also proves the regularity needed to return
to admissible paths.  Put $A=\alpha'$ and let $Y$ be the variation field.
The derivative $D_s$ is the pullback covariant derivative for $g(T-s^2)$.
Metric compatibility on each slice and $\partial_sg(T-s^2)=4s\operatorname{Ric}$
give
\[
 \frac d{ds}\langle A,Y\rangle
 =\langle D_sA,Y\rangle+\langle A,D_sY\rangle
       +4s\operatorname{Ric}(A,Y).
\]
Thus differentiating \eqref{eq:regularized-action} gives
\begin{equation}\label{eq:first-variation}
 \delta\mathcal L=[\langle A,Y\rangle]_{s_1}^{s_2}
 -\int_{s_1}^{s_2}
 \left(\langle D_sA,Y\rangle-2s^2dR(Y)
                       +4s\operatorname{Ric}(A,Y)\right)ds.
\end{equation}
Compactly supported variations give the weak Euler equation.  In coordinates
it first regularizes the momentum, then the velocity; bootstrapping the
smooth coefficients gives a smooth solution.  Its intrinsic equations are
\[
 D_sA-2s^2\nabla R+4s\operatorname{Ric}^{\sharp}(A)=0,
 \qquad
 D_\tau X-\tfrac12\nabla R+\frac X{2\tau}
                  +2\operatorname{Ric}^{\sharp}(X)=0,
 \quad X=\dot\gamma.
\]
The first is nonsingular at $s=0$.  Coordinate ODE uniqueness identifies
overlapping local solutions.  Endpoint continuation uses bounds for the
actual phase curve in a compact chart neighborhood, not merely the formal
absence of a singular coefficient.  The resulting square-root representative
is smooth near the closed interval and hence gives the required minimizer.
The same phase estimates and ODE uniqueness continue a positive-time Euler
path down to zero.

Here is the index form, including the terms caused by the evolving metric.
Let $\mathrm{Rm}$ denote the four-slot curvature tensor with the slot convention
of the curvature chapter.  For a square-root variation of an Euler path,
the second derivative is $[\langle A,Z\rangle]_{s_1}^{s_2}+I(Y,Y)$,
where $Z$ is the covariant acceleration of its endpoint variation and
\begin{align*}
 I(Y,Y)=\int_{s_1}^{s_2}\bigl(&|D_sY|^2+\mathrm{Rm}(Y,A,A,Y)
       +2s^2\nabla^2R(Y,Y)\\
       &-4s(\nabla_Y\operatorname{Ric})(A,Y)
             +2s(\nabla_A\operatorname{Ric})(Y,Y)\bigr)\,ds.
\end{align*}
To obtain the last two terms, differentiate the connection in the family:
its backward-time variation pairs with $U,V,W$ as
$(\nabla_U\operatorname{Ric})(V,W)+
 (\nabla_V\operatorname{Ric})(U,W)-
 (\nabla_W\operatorname{Ric})(U,V)$.  Commuting variation and pullback
derivatives and integrating by parts gives the displayed formula.
This also defines the Jacobi operator by linearizing the Euler equation.
At a minimizer $I$ is nonnegative on fields vanishing at both endpoints.
If $I(Y,Y)=0$, positivity of $I(Y+tZ,Y+tZ)$ for every real $t$ gives
$I(Y,Z)=0$ for every such $Z$; integration by parts and the fundamental
lemma give the Jacobi equation.  Conversely a Jacobi field with both endpoint
values zero has zero index.  Its initial-zero, prescribed-initial-derivative
problem is a nonsingular linear ODE in $s$, so it has a unique solution;
the prescribed derivative does not impose a terminal-zero condition.
\end{proof}

\section{Differential geometry of reduced length}

The minimizing curve is not unique everywhere.  The correct regular set is
therefore a set of endpoint--time pairs, not a globally chosen branch.

\begin{definition}[Regular branch and Harnack integral]
\label{def:reduced-regular-branch}
\lean{PoincareMT.ReducedLengthRegularPoint,PoincareMT.reducedHarnackIntegral}
\uses{def:reduced-length}
For a base point $p$, a point $(q,\tau)$ with $0<\tau<\tau_{\max}$ is regular
if it has a unique minimizing path whose endpoint map has bijective
differential.  Write $l(q,\tau)=\ell_T(p,q,\tau)$ also at nonregular points.
On a regular branch, if
$R_\tau(\sigma)$ is the derivative of
$r\mapsto R(\gamma(\sigma),T-r)$ at $r=\sigma$, with its spatial point
held fixed, set
\[
 H(\tau)=-R_\tau(\tau)-\frac{R(\tau)}\tau
 -2\langle\nabla R,\dot\gamma\rangle+2\operatorname{Ric}(\dot\gamma,\dot\gamma),
 \quad K(\tau)=\int_0^\tau s\sqrt{s}\,H(s)\,ds .
\]
In particular $R_\tau$ is not the total derivative along $\gamma$;
the separate spatial term in $H$ must be retained.
\end{definition}

\begin{theorem}[Regular formulas and barriers]
\label{thm:m09-differential}
\uses{thm:m08-l-geodesics,def:reduced-regular-branch,prop:rf-evolution}
\lean{PoincareMT.ReducedLengthDifferentialTheory}
\lean{PoincareMT.reducedLengthDifferentialInequalities}
Under the complete bounded-curvature hypotheses above, for every $p$ and
every $0<\tau<\tau_{\max}$ there is a dense nonempty open
set of regular endpoints.  On a regular branch,
\begin{align*}
 \partial_\tau l&=R-\frac l\tau+\frac K{2\tau\sqrt\tau},\\
 |\nabla l|^2&=\frac l\tau-\frac K{\tau\sqrt\tau}-R,\\
 \Delta l&\leq\frac n{2\tau}-R-\frac K{2\tau\sqrt\tau},\\
 \partial_\tau l+\Delta l&\leq\frac{n/2-l}{\tau},\\
 \partial_\tau l-\Delta l+|\nabla l|^2-R+\frac n{2\tau}&\geq0,\\
 2\Delta l-|\nabla l|^2+R+\frac{l-n}{\tau}&\leq0.
\end{align*}
At every endpoint, and every $\varepsilon>0$, there is a smooth upper barrier
$b\geq l$ touching $l$ at that endpoint with residual
\[
 \partial_\tau b+\Delta b-\frac{n/2-b}{\tau}\leq\varepsilon .
\]
For every positive-time spacetime point there are a neighborhood $N$ and a
constant $C\geq0$ such that every point of $N$ has an upper barrier with
$|\partial_\tau b|\leq C$, $|\nabla b|^2\leq C$, and
$\nabla^2b(v,v)\leq Cg(v,v)$ there at its touching point.  This is an
upper Hessian bound, not a two-sided Hessian norm bound.  On a
regular branch equality in the sharp Laplacian inequality implies
\[
 \operatorname{Ric}+\nabla^2l=\frac{g}{2\tau}.
\]
\source{MT2007}\dcref{Propositions 6.37 and 6.43, pp. 123--129}
\dcref{Theorem 6.50 and Corollary 6.51, pp. 131--132}
\dcref{Propositions 6.56 and 6.59, pp. 133--139}
\dcref{Proposition 7.8, pp. 152--154}\cite[Propositions 6.37 and 6.43, Theorem 6.50, Corollary 6.51, Propositions 6.56 and 6.59, and Proposition 7.8]{MT2007}
\end{theorem}
\begin{proof}
The coherent $\mathcal L$-exponential family is smooth for every initial
vector and every $0<\tau<\tau_{\max}$, including nonminimizing branches.
Indeed, on each strict subslab the curvature derivative estimate bounds
$|\nabla R|$; the square-root Euler equation then bounds speed and
displacement.  Completeness confines the phase to a compact set, where
finitely many ODE neighborhoods provide a common restart interval.
Maximality therefore extends each solution through that subslab, and smooth
ODE dependence gives the family.  Its minimizing inverse branches remain
smooth until the differential loses invertibility or a competitor appears.
The complement of these events is open; strict prefixes of minimizing
branches remain minimizing, and the cut-time and critical-value arguments
give density.  On a regular inverse branch, \eqref{eq:first-variation} gives
$\nabla l=X(\tau)$ and endpoint differentiation gives
$\partial_\tau l=(R-|X|^2)/2-l/(2\tau)$, with the endpoint held fixed.
Integrating the derivative of $\tau^{3/2}(R+|X|^2)$ along an Euler path,
using scalar evolution, gives
$|X|^2=l/\tau-R-K/(\tau\sqrt\tau)$.
These two equations give the first two displayed identities.  Apply the
index form to endpoint-adapted comparison fields, sum over an orthonormal
endpoint basis, and use the contracted Bianchi identity in the Ricci-derivative
terms.  This gives the displayed Laplacian estimate.  Adding it to the
time identity gives the fourth inequality; subtracting it and substituting
the gradient identity gives the fifth; doubling it and subtracting the
gradient identity gives the sixth.  Thus all coefficients of $K$ cancel
in the three inequalities.

At a singular endpoint choose a minimizing path and write $c=\sqrt\tau$.
Transport an orthonormal frame $P_i(s)$ along its square-root representative
by $D_sP_i=-2s\operatorname{Ric}^{\sharp}P_i$.  The moving-metric identity
shows that the frame remains orthonormal.  Realize the fields
$Y_i(s)=(s/c)P_i(s)$ simultaneously by a smooth finite-parameter family
$f(z,s)$ with $f(0,s)=\alpha(s)$ and $f(z,0)=p$.
Such a family is obtained from local smooth variations along the compact
curve and a partition of unity.  Its terminal differential sends the
coordinate basis to $P_i(c)$ and is invertible.  Thus
$(z,\sigma)\mapsto(f(z,\sqrt\sigma),\sigma)$ has a smooth local inverse,
even when the minimizing exponential map is critical.

Define $b(q,\sigma)$ to be the action of this comparison path divided by
$2\sqrt\sigma$.  Its central action is $2\sqrt\tau\,l$, and every nearby
path is an admissible competitor, so $b$ touches $l$ from above.
First variation gives the same time derivative as for the central minimizer;
second variation in the $n$ parameter directions gives
\[
 \Delta b=\frac n{2\tau}-R-\frac K{2\tau\sqrt\tau}
\]
at contact.  Endpoint acceleration terms cancel against the connection
terms in the Hessian of $b$.  The Harnack identity now gives
$\partial_\tau b+\Delta b-(n/2-b)/\tau=0$, which implies the asserted
$\varepsilon$ bound.  Compact chart and minimizing-speed bounds give the
stated locally uniform first-derivative and upper-Hessian bounds.

For the regular equality assertion, let $J_i$ be the Jacobi field with the
same two endpoint values as $Y_i$.  Integration by parts gives
$I(Y_i,Y_i)-I(J_i,J_i)=I(Y_i-J_i,Y_i-J_i)\geq0$.
Equality in the trace forces each difference to have zero index.  It is
then an endpoint-zero Jacobi field, hence zero by noncriticality.  Thus the
Jacobi fields satisfy the adapted equation.  Torsion freeness and
$X=\nabla l$ give $D_\tau J_i=\nabla_{J_i}\nabla l$ at the endpoint;
the adapted equation gives
$D_\tau J_i=-\operatorname{Ric}^{\sharp}J_i+J_i/(2\tau)$.
Since their endpoint values form a basis, the claimed tensor identity follows.
\end{proof}

\section{Reduced volume and rigidity}

\begin{definition}[Unnormalized reduced volume]
\label{def:reduced-volume}
\lean{PoincareMT.reducedVolume,PoincareMT.reducedVolumeDensity}
\uses{def:reduced-length}
At positive backward time define
\[
 u_T(p,q,\tau)=\tau^{-n/2}e^{-\ell_T(p,q,\tau)},\qquad
 V_{\mathrm{red},T}(p,\tau)=\int_Mu_T(p,q,\tau)\,d\mu_{g(T-\tau)}(q).
\]
The density is extended by zero at $\tau\leq0$ and uses calibrated Riemannian
Hausdorff measure on the selected slice.  With this convention Euclidean space
has
\[
 V_{\mathrm{red},\mathbb R^n}(p,\tau)=(4\pi)^{n/2};
\]
there is no hidden $(4\pi)^{-n/2}$ normalization.
\end{definition}

\begin{theorem}[Weak inequalities, monotonicity and equality]
\label{thm:m10-reduced-volume}
\uses{def:reduced-volume,thm:m09-differential}
\lean{PoincareMT.reducedVolumeMonotonicity}
Under the hypotheses of Theorems~\ref{thm:m08-l-geodesics} and
\ref{thm:m09-differential}, $l$ is locally Lipschitz in positive spacetime,
smooth on its open regular locus, whose complement has zero measure on
each positive-time slice.  It satisfies the weak inequalities
\eqref{eq:reduced-weak-inequalities}.  At every $0<\tau<\tau_{\max}$, $u_T$ is
integrable, $V_{\mathrm{red},T}(p,\tau)>0$ and
\(V_{\mathrm{red},T}(p,\tau)\leq(4\pi)^{n/2}\).  It tends to $(4\pi)^{n/2}$ as
$\tau\downarrow0$ and is nonincreasing on $(0,\tau_{\max})$.  The same
monotonicity holds after restricting to any open regular domain
$A\subset M\times(0,\tau_{\max})$ such that, for each $(q,\tau)\in A$,
its minimizing branch satisfies $(\gamma(\sigma),\sigma)\in A$ for every
$0<\sigma\leq\tau$.  Also $l(\cdot,\tau)$ attains a spatial minimum
at most $n/2$.  If equality with the Euclidean value holds at an interior
time $b$, the closed slab $[T-b,T]$ is static Euclidean: one fixed
diffeomorphism identifies every metric on that slab with the Euclidean metric.
\source{MT2007}\dcref{Corollary 6.67 and Claims 6.68--6.69, pp. 139--140}
\dcref{Definition 6.70 and Lemma 6.71, pp. 140--141}
\dcref{Theorem 6.80 and Proposition 6.81, pp. 145--146}
\dcref{Theorem 7.10 and Claim 7.11, pp. 154--156}
\dcref{Theorem 7.13, pp. 156--164}
\dcref{Claim 7.22 and Corollary 7.24, pp. 161--162}
\dcref{Theorem 7.26 and Proposition 7.27, pp. 165--167}\cite[Corollary 6.67, Claims 6.68--6.69, Lemma 6.71, Theorem 6.80, Proposition 6.81, Theorems 7.10 and 7.13, Claim 7.22, Corollary 7.24, Theorem 7.26, and Proposition 7.27]{MT2007}
\end{theorem}
\begin{proof}
Fix $0<\tau<\tau_{\max}$; all integrals in the weak-comparison argument are over
$(M,g(T-\tau))$.  For a nonnegative compactly supported smooth test function $\varphi$, the
regular formulas and the barrier extension give the two weak inequalities
\begin{equation}\label{eq:reduced-weak-inequalities}
\begin{aligned}
 \int_M\!\left[\varphi\left(\partial_\tau l+|\nabla l|^2-R+\frac n{2\tau}\right)
 -l\Delta\varphi\right]d\mu&\geq0,\\
 \int_M\!\left[\varphi\left(-|\nabla l|^2+R+\frac{l-n}{\tau}\right)
 +2l\Delta\varphi\right]d\mu&\leq0.
\end{aligned}
\end{equation}
Here is why the nonregular set does not obstruct this step.  The local
upper-barrier bounds imply local Lipschitz regularity of $l$.  At almost
every endpoint its spatial differential exists, by Rademacher's theorem in
charts.  The critical values of the smooth exponential map have zero measure
by Sard's theorem.  At a differentiable endpoint outside those critical
values, the first variation identifies the terminal covector of every
minimizer with $dl$.  Thus two minimizing lifts have the same terminal
position and velocity.  Uniqueness for the coordinate phase ODE identifies
them near that endpoint; uniqueness of extension to zero identifies their
whole paths.  To verify neighborhood stability, suppose instead that
minimizing lifts to endpoints approaching $q$ leave a fixed inverse-function
neighborhood of its unique lift.  Local action and speed bounds make their
initial vectors bounded.  A subsequence converges to a minimizing lift of
$q$, by continuous dependence and continuity of the action and reduced length.
Uniqueness identifies this limit with the original lift, contradicting that
the sequence stayed outside its neighborhood.  This gives a regular branch.
The slice complement of the regular
locus is consequently null.  Positive smooth chart volume densities transfer
these null sets to the calibrated slice measure.

An almost-everywhere Laplacian inequality alone would not control a singular
part of the distributional Hessian.  The upper-barrier Hessian bounds supply
the missing sign.  In a smaller chart, after accounting for the bounded
Christoffel term with the gradient bound, there is $K$ such that
$l(y)-K|y|^2/2$ is concave.  Its distributional Hessian is a negative
semidefinite matrix of measures.  If $B(y)$ is the coordinate metric and
$\rho(y)=\sqrt{\det B(y)}$, the coordinate Laplacian is
\[
 \rho\Delta f=\operatorname{div}(\rho B^{-1}df).
\]
The singular Hessian contribution has nonpositive contraction with the
positive matrix $\rho B^{-1}$.  Concretely, smooth convolution on an inner
chart gives $l_j\to l$ uniformly, a common upper Hessian bound and a common
gradient bound.  The bounds on $B,B^{-1},\rho$ and their first derivatives
give $\Delta l_j\leq C$ on the compact test support.  At every regular
point the derivatives through order two converge, so
$\Delta l_j\to\Delta l$ almost everywhere.  Apply Fatou's lemma to
$\varphi(C-\Delta l_j)\geq0$ and integrate by parts against the fixed
smooth test.  This carries the almost-everywhere upper bound $\Delta l\leq h$ to
$\int l\Delta\varphi\leq\int\varphi h$ for nonnegative tests supported
there.  Here $h=-\partial_\tau l+(n/2-l)/\tau$ is locally bounded and
measurable.  The regular formulas identify the first displayed integrand
almost everywhere with $\varphi h-l\Delta\varphi$, and the second with
minus twice that expression.  A finite smooth partition of unity on the
test support gives both global inequalities and their integrability.
No pointwise derivative on the cut locus is required.

For the minimum assertion, the quadratic action lower bound makes every
bounded sublevel of $l(\cdot,\tau)$ compact.  Continuity gives an attained
minimum $m(\tau)$, continuous on positive compact time intervals.
At a minimizing point $q$ an upper barrier also has a spatial local minimum,
so $\Delta b(q,\tau)\geq0$.  Its residual inequality implies
\[
 b(q,\tau)+\tau\partial_\tau b(q,\tau)
       \leq n/2+\tau\varepsilon.
\]
The function $s\mapsto s b(q,s)$ is an upper support for $s m(s)$.
Letting $\varepsilon\downarrow0$ and applying one-sided slope comparison
gives $b m(b)-a m(a)\leq(n/2)(b-a)$ for $0<a\leq b<\tau_{\max}$.
The scalar lower bound and a short comparison path at the base show
$m(a)\to0$ as $a\downarrow0$.  Thus $m(b)\leq n/2$.

For monotonicity the construction uses a change of variables to the initial
tangent space.  Normalize its coordinate $Z$ by
$\lim_{\tau\downarrow0}\sqrt\tau\,\dot\gamma_Z(\tau)=Z$.
Thus in Euclidean space $\gamma_Z(\tau)=p+2\sqrt\tau Z$.
Let $D_\tau$ be the open set of regular minimizing initial vectors and
$J(Z,\tau)$ the positive Jacobian of $Z\mapsto\gamma_Z(\tau)$ relative
to $g(T)$ on the source and $g(T-\tau)$ on the target.  Set
\[
 w_\tau(Z)=\boldsymbol 1_{D_\tau}(Z)\,
       \tau^{-n/2}e^{-l(\gamma_Z(\tau),\tau)}J(Z,\tau).
\]
On a regular ray, the slice-volume derivative and $X=\nabla l$ give
$\partial_\tau\log J=R+\Delta l$.  The total derivative of the endpoint
action is $d_\tau l(\gamma_Z(\tau),\tau)=\partial_\tau l+|\nabla l|^2$.
Consequently, with
$\delta=n/(2\tau)-R-K/(2\tau\sqrt\tau)-\Delta l\geq0$,
\[
 \partial_\tau\log(\tau^{-n/2}e^{-l(\gamma_Z(\tau),\tau)}J)
 =-\delta.
\]
This proves that the weight is nonincreasing along every regular minimizing
ray.  The strict-prefix
property gives $D_b\subset D_a$ for $0<a\leq b$, so extension by zero
preserves the comparison even after a ray leaves the regular domain.
The small-time comparison argument developed in the proof of
Theorem~\ref{thm:m14-generalized}, applied to the ordinary product, puts
each fixed compact set of initial vectors in $D_\tau$ for all sufficiently
small positive $\tau$.  That argument uses only local ODE and action
estimates, without a volume conclusion.  The initial Jacobi derivative and action expansion now give
\[
 \lim_{\tau\downarrow0}w_\tau(Z)=2^ne^{-|Z|^2},\qquad
 0\leq w_\tau(Z)\leq 2^ne^{-|Z|^2}.
\]
The complement of the regular image has zero slice measure.  Calibrated
change of variables therefore gives
$V_{\mathrm{red},T}(p,\tau)=\int_{T_pM}w_\tau(Z)\,dZ$.
The Gaussian is integrable, with integral $(4\pi)^{n/2}$; integral comparison
gives monotonicity and the upper bound, and dominated convergence gives the
zero-time limit.  Positivity follows from a nonempty open regular image.
For the restricted domain $A$, insert the additional indicator
$\boldsymbol1_A(\gamma_Z(\tau),\tau)$: its prefix property gives the same
pointwise comparison.  This argument does not require a global smooth
reduced-length function on the cut locus.

If $V_{\mathrm{red},T}(p,b)=(4\pi)^{n/2}$, the nonnegative difference
$2^ne^{-|Z|^2}-w_b(Z)$ has integral zero.  Equality propagates to all
$0<s<b$ by monotonicity.  To see why the exponential becomes global and
one-to-one, the equality first makes its regular source conull, hence dense.
Here the raw weight uses the action of the full smooth exponential family,
without the indicator of $D_s$; that action equals the minimizing action on
$D_s$.  The raw weighted Jacobian is continuous; its almost-everywhere equality with
the strictly positive Gaussian therefore holds everywhere.  Its positive
Gram determinant makes the differential invertible at every initial vector,
so the inverse function theorem makes the full exponential map open.
It is injective on its open dense regular source.  An open continuous map
with this property is globally injective: if disjoint neighborhoods of two
points mapped over the same point, density first gives a regular source
point in one neighborhood mapping into the image of the other; openness and
density then give a regular point in the other mapping into the first
regular neighborhood.  This contradicts injectivity on the regular source.
Every endpoint has a minimizing lift by existence, so the full map is also
surjective.  Its smooth local inverses consequently glue to a global smooth
inverse.  The logarithmic derivative just computed is zero, so $\delta=0$.
The regular equality assertion yields
$\operatorname{Ric}+\nabla^2l=g/(2s)$ at every endpoint for $0<s<b$.
Here $s$ denotes backward time and $\gamma_s$ is the map
$Z\mapsto\gamma_Z(s)$.  For a fixed source vector $U$, let $Y(s)=d\gamma_sU$ and
$a(s)=g(T-s)(Y(s),Y(s))$.  Differentiating, using
$D_sY=\nabla_Y X$ and $X=\nabla l$, gives
\[
 a'(s)=2(\operatorname{Ric}+\nabla^2l)(Y,Y)=a(s)/s.
\]
Thus $a(s)/s$ is constant.  The initial differential
$d\gamma_sU/(2\sqrt s)\to U$ fixes its value as $4g(T)(U,U)$.
Polarization now gives the pairing identity
\[
 g(T-s)(d\gamma_s U,d\gamma_s W)=4s\,g(T)(U,W).
\]
The determinant of this identity gives $J(Z,s)=(2\sqrt s)^n$.
Substitution in the Gaussian weight equality gives
$l(\gamma_Z(s),s)=|Z|_{g(T)}^2$.  Differentiate this identity in a source
direction $U$.  First variation gives
$g(T-s)(\partial_s\gamma_Z(s),d\gamma_sU)=2g(T)(Z,U)$.
The pairing identity gives exactly the same pairing for
$d\gamma_s(Z/(2s))$.  Surjectivity of $d\gamma_s$ and positive definiteness
therefore imply the radial-velocity identity
$\partial_s\gamma_Z(s)=d\gamma_s(Z/(2s))$.
Consequently the rescaled map
$f_s(y)=\gamma_{y/(2\sqrt s)}(s)$ has zero derivative in $s$:
the derivative of its initial vector is $-Z/(2s)$ and cancels the radial
velocity.  Fix $f=f_{b/2}$.  The pairing identity says $f^*g(T-s)$ is the
fixed Euclidean metric for all $0<s<b$.  Continuity gives the same identity
at $s=0,b$.  It is this stationary map, not merely a family of slice
isometries, that proves the claimed rigidity.
\end{proof}

\section{The generalized spacetime carrier}

An ordinary product $I\times M$ is insufficient after singular limits.  The
generalized construction keeps a single topological carrier $X$ and a smooth
time function $t:X\to\mathbb R$ whose range is the prescribed interval.

\begin{definition}[Generalized flow spacetime]
\label{def:generalized-spacetime}
\lean{PoincareMT.GeneralizedFlowSpacetime}
A generalized flow spacetime consists of a Hausdorff second-countable
charted smooth carrier with half-space time model, a smooth clock with
range $I$ and $\partial X=\{x:t(x)\in\operatorname{frontier}I\}$,
a smooth normalized time vector $\chi$
with $dt(\chi)=1$, and the horizontal bundle
\[
 \mathcal H_x=\ker(dt_x)\subset T_xX.
\]
It includes a smooth positive metric on $\mathcal H$ and a smooth projection
\[
 \Pi_x(v)=v-dt_x(v)\chi_x\in\mathcal H_x,
 \qquad v=dt_x(v)\chi_x+\Pi_x(v).
\]
The slice at time $s$ is the inherited level set $\{x:t(x)=s\}$; its charts
and metric are separate data obtained from the carrier.
\end{definition}

An adapted atlas has time-preserving product charts; on overlaps the spatial
change of coordinates preserves the supplied metric and time direction.
A compatible cylinder is a smooth embedding $F:K\times N\to X$ satisfying
$t(F(s,q))=s$ and $dF(\partial_s)=\chi$, with the pulled-back horizontal
metric as its spatial metric.  A gauge cover is a family of such cylinders
covering the carrier, together with their local inverses.  It asserts no
Ricci equation.  For horizontal fields $U,V$, define
\[
 (\mathcal L_\chi g)(U,V)
 =\chi(g(U,V))-g(\Pi[\chi,U],V)-g(U,\Pi[\chi,V]).
\]
The intrinsic Ricci-flow equation is
$\mathcal L_\chi g=-2\operatorname{Ric}^{\mathcal H}$.

\begin{theorem}[Carrier, gauge and Ricci bridge]
\label{thm:m11-m12-gauge}
\uses{def:generalized-spacetime}
\lean{PoincareMT.generalizedSpacetimeRealization}
\lean{PoincareMT.GeneralizedRicciGaugeTheory}
\lean{PoincareMT.generalizedRicciGaugeGeometry}
For every adapted metric atlas, the construction supplies such a carrier and
horizontal metric on the exact original topology.  For every supplied gauge
cover, there is a leafwise Levi--Civita family and horizontal Ricci calculus.
In each gauge chart the intrinsic horizontal Ricci equation is equivalent to
the ordinary Ricci PDE for the chart metric.  In the ordinary product case the
same cylinder metric and slice metrics are used by both sides, and an intrinsic
solution yields an ordinary-flow gauge witness.  The selected interval
structures also have smooth restriction maps: an interval inclusion whose
image is relatively open is a local diffeomorphism.
\source{MT2007}\dcref{Definitions 3.34--3.36 and Remark 3.37, pp. 59--60}\cite[Definitions 3.34--3.36 and Remark 3.37, pp. 59--60]{MT2007}
\end{theorem}
\begin{proof}
The adapted atlas constructs a charted manifold on the given topological
space; the common chart clock gives its time function and time range.
The horizontal bundle is literally
the kernel of $dt$, rather than an abstract copy of a spatial tangent bundle;
the projection formula above proves the direct-sum decomposition and its
smoothness.  In a product chart this is the usual splitting into time and
spatial coordinates, so the spatial transition differentials give local
trivializations of the kernel bundle.  Metric compatibility on overlaps
glues the positive horizontal forms.  Level-set charts give the slices and
identify their tangent spaces with these horizontal fibers.

Pull the horizontal metric to each gauge leaf.  In coordinates its
Levi--Civita coefficients are
$\Gamma^k_{ij}=\tfrac12g^{ka}(\partial_i g_{ja}+\partial_j g_{ia}
-\partial_a g_{ij})$.  Smooth inversion of the positive metric matrix
gives smooth dependence on time.  On overlapping gauges, uniqueness of the
torsion-free metric connection identifies these connections, hence also
their curvature and Ricci contractions.  For coordinate fields transported
by the cylinder, $[\chi,U]=[\chi,V]=0$.  The Lie derivative above therefore
pulls back to $\partial_tg_{ij}$, and the intrinsic equation pulls back
to the ordinary Ricci PDE.  Conversely, every point and horizontal pair
lifts to a gauge, so the ordinary equations imply the tensor equation.
Levi--Civita uniqueness makes this equivalence independent of the selected
connection representative.  In the product construction the identifications
are the given $(s,q)$ and $q$ themselves.  For a relatively open interval
restriction, the inverse is the identity on the same time coordinate;
restriction of its interval charts proves the local-diffeomorphism assertion.
No completeness follows from these local constructions.
\end{proof}

\section{Rescaling and generalized reduced geometry}

\begin{theorem}[Exact parabolic rescaling]
\label{thm:m13-rescaling}
\uses{thm:m11-m12-gauge,def:reduced-length,def:reduced-volume}
\lean{PoincareMT.ParabolicSpacetimeRescaling}
\lean{PoincareMT.generalizedParabolicRescaling}
Given $Q>0$ and a time shift $a$, the target has clock $t'=Q(t-a)$, metric
$g'=Qg$, time vector $\chi'=\chi/Q$, and the precisely transported interval.
The identity on the carrier is a diffeomorphism; its differential is the
identity on tangent vectors, it identifies horizontal fibers, and the slice
map is a metric homothety of factor $Q$.  Connections are unchanged under
this constant homothety.  Scalar curvature and the full curvature norm scale
by $Q^{-1}$, while
the Ricci tensor as a covariant two-tensor is unchanged.  Volume has factor
$Q^{n/2}$.  A normalized initial tangent is divided by
$\sqrt Q$; hence $\mathcal L$ scales by $\sqrt Q$ while reduced length is
unchanged on its finite domain.
\source{MT2007}\dcref{Definition 3.40, p. 61}\cite[Definition 3.40, p. 61]{MT2007}
\end{theorem}
\begin{proof}
The transformed time interval is $\{Q(t-a):t\in I\}$, with the same
endpoint-inclusion conventions as $I$; the inverse clock is $a+t'/Q$.
The time differential is $dt'=Qdt$, so the normalized vector must be
$\chi/Q$.  The horizontal kernel is unchanged and the identity differential
therefore gives the fiber equivalence.  Constant metric homothety leaves the
Levi--Civita connection unchanged and multiplies the quadratic speed and
volume density by $Q$ and $Q^{n/2}$ respectively.  For the reparameterized
path $\gamma'(\tau')=\gamma(\tau'/Q)$, the squared speed and scalar
curvature both acquire the factor $Q^{-1}$.  Hence
$\sqrt{\tau'}(R'+|\dot\gamma'|^2)d\tau'
=\sqrt Q\sqrt\tau(R+|\dot\gamma|^2)d\tau$.
The action has factor $\sqrt Q$; the denominator
$2\sqrt{\tau'}=2\sqrt Q\sqrt\tau$ cancels it in reduced length.  The square-root
initial derivative forces the $1/\sqrt Q$ tangent normalization.  These are
equalities of the selected carrier data, not merely abstract isomorphisms.
\end{proof}

From now on the generalized carrier is supplied with a gauge cover, its
leafwise Levi--Civita family, and the intrinsic equation
$\mathcal L_\chi g=-2\operatorname{Ric}^{\mathcal H}$.
The gauge bridge therefore makes each compatible cylinder an ordinary
Ricci flow with exactly its pulled-back metric.

\begin{definition}[Generalized backward paths and stable vectors]
\label{def:generalized-stable-paths}
\lean{PoincareMT.M14BackwardPath,PoincareMT.M14StableSet}
\uses{thm:m11-m12-gauge,def:reduced-regular-branch}
A backward path based at $(x,T)$ is a curve $\gamma$ in $X$ with
$t(\gamma(\tau))=T-\tau$, prescribed endpoints, continuous closed-interval
and $C^1$ interior regularity, and integrable action.  Its spatial velocity is
$v=\Pi\dot\gamma$; the clock equation gives $\dot\gamma=-\chi+v$.
Its action is
$\int\sqrt\tau(R^{\mathcal H}+g(v,v))\,d\tau$.
The action infimum is used as a finite real value only when its competitor
set is nonempty and bounded below.

An exponential family $E(Z,s)$ is the coherent square-root Euler solution
with $E(Z,0)=x$ and initial horizontal velocity $2Z$.  Its domain is
relatively open in $\{(Z,s):s\geq0,\ T-s^2\in I\}$ and is exactly
the survival domain of those initial-value solutions.  At fixed $\tau>0$,
a vector is stable if it survives to $\sqrt\tau$, the endpoint differential
is bijective, and it has an open neighborhood all of whose branches uniquely
minimize to their respective endpoints.  The stable set consists of exactly
these vectors.  It has a smooth endpoint map into $X_{T-\tau}$ and smooth
local inverses.  It may be empty.  A pointwise uniquely minimizing branch
alone is not the full neighborhood stability condition.
For a measurable set $W$ of stable vectors define
\[
 V_{\mathrm{red}}(W,\tau)=
 \int_{E(W,\sqrt\tau)}\tau^{-n/2}e^{-l(q,\tau)}\,d\mu_{T-\tau}(q).
\]
Unique minimizing branches make the endpoint map injective there, so this
also equals the integral over $W$ of its weighted Jacobian.  The generalized
zero-time limit below concerns this volume on the full stable image;
it makes no assertion about the complement of that image in the slice.
\end{definition}

\begin{proposition}[Positive starting time]
\label{prop:positive-start}
\uses{def:generalized-stable-paths,thm:m08-l-geodesics}
\lean{PoincareMT.M14.positiveStartCorrectionStatement}
Let $0<a<b$, and let $\gamma:[a,b]\to X$ be a minimizing backward
path from $x$ to $y$, with its genuine one-sided initial derivative
$\dot\gamma(a)=-\chi+v(a)$.  Suppose the endpoint reduced-length function
$l(q)=\mathcal L_{\min}(x,q)/(2\sqrt{T-t(q)})$ is smooth on an open
neighborhood of $y$ contained in $\{q:a<T-t(q)\}$.  All differential
operators below are on the selected horizontal slice, and
$\partial_\tau l=-\chi(l)$.  Define $H$ as above using horizontal
curvature and velocity, and put
\[
 C_a=(a/b)^{3/2}(R(\gamma(a))+|v(a)|^2),\quad
 K_a=\int_a^b s^{3/2}H(s)\,ds,\quad
 K_a^{\mathrm{shift}}=\int_a^b\sqrt s(\sqrt s-\sqrt a)^2H(s)\,ds.
\]
Then, with $l=l(y)$ and $R=R(y)$,
\begin{align*}
 \partial_\tau l&=R-l/b+K_a/(2b\sqrt b)-C_a/2,\\
 |\nabla l|^2&=l/b-K_a/(b\sqrt b)-R+C_a,\\
 \Delta l&\leq\frac{n}{2\sqrt b(\sqrt b-\sqrt a)}-R
       -\frac{K_a^{\mathrm{shift}}}
                    {2\sqrt b(\sqrt b-\sqrt a)^2}.
\end{align*}
\source{MT2007}\dcref{Proposition 6.43, pp. 127--129}
\dcref{Lemma 6.49 and Theorem 6.50, p. 131}\cite[Proposition 6.43, pp. 127--129; Lemma 6.49 and Theorem 6.50, p. 131]{MT2007}
\end{proposition}
\begin{proof}
The first variation identifies the endpoint gradient with $v(b)$; moving
the terminal time with spatial endpoint fixed gives
$\partial_\tau l=(R-|v(b)|^2)/2-l/(2b)$.
Integrating the Euler/Harnack identity from $a$ to $b$ leaves its lower
endpoint value $a^{3/2}(R(\gamma(a))+|v(a)|^2)$.
Dividing by $b^{3/2}$ produces $C_a$ and gives the gradient identity;
substitution gives the time identity.  For the Laplacian comparison, use
endpoint-adapted fields vanishing at $\sqrt a$, with linear factor
$(s-\sqrt a)/(\sqrt b-\sqrt a)$ in square-root time.  The leading
index term gives $n/(2\sqrt b(\sqrt b-\sqrt a))$, and the Harnack term
is weighted by the square of that factor, giving $K_a^{\mathrm{shift}}$.
Keeping the original clock is essential: replacing $b$ by $b-a$ in the
zero-start formulas would discard the initial term and change the action.
\end{proof}

\begin{theorem}[Generalized $\mathcal L$-geometry and transport]
\label{thm:m14-generalized}
\uses{prop:positive-start,thm:m10-reduced-volume,thm:m13-rescaling}
\lean{PoincareMT.generalizedLGeometryTheory}
Assume the curvature calculus of the Ricci-flow chapter and the gauge and
rescaling constructions above.
On every generalized Ricci flow with the supplied gauge and slice data,
actual backward paths satisfy the clock equation
and use horizontal velocity.  On their genuine domains the square-root Euler
equation, first and second variation, Jacobi IVP, stable minimizing sets,
regular formulas, calibrated Jacobian transport and reduced-volume transport
hold.  Finite action alone does not assert attainment.  If an ordinary
complete bounded-curvature presentation is supplied, the same carrier
recovers the ordinary existence, differential and volume conclusions through
their ordinary-flow theorems.  Suppose $\delta>0$,
$[T-\delta,T]\subset I$, and one constant bounds horizontal full curvature
at every point of that entire past slab.  For every compact set $B$ of initial
horizontal vectors and every supplied compact neighborhood of the base point,
there is $0<\varepsilon\leq\delta$ such that $B$ lies in the stable domain
at every $0<\tau\leq\varepsilon$ and its branches stay in that neighborhood.
With compact exhaustion the reduced volume tends to $(4\pi)^{n/2}$ at zero.
\source{MT2007}\dcref{Definitions 6.1--6.2, p. 106}
\dcref{Definition 6.25 and Proposition 6.28, pp. 116--118}
\dcref{Proposition 6.78 and Corollary 6.79, pp. 143--145}
\dcref{Theorem 6.80 and Proposition 6.81, pp. 145--146}
\dcref{Proposition 7.5, pp. 151--152}
\dcref{Theorem 7.26, p. 165}\cite[Definitions 6.1--6.2 and 6.25, Proposition 6.28, Proposition 6.78, Corollary 6.79, Theorem 6.80, Proposition 6.81, Proposition 7.5, and Theorem 7.26]{MT2007}
\end{theorem}
\begin{proof}
The generalized path calculus is proved in gauge cylinders and transported
between overlapping cylinders by the carrier's horizontal projection.  The
square-root change of variable removes the initial singular coefficient, as
in the ordinary variational calculation.  Stable sets are defined by unique minimizing branches and are closed
under strict prefixes; this gives a single exponential family rather than a
choice made independently at each endpoint.  The Jacobian of that family is
computed with the Jacobi fields.  Multiplying the Jacobian by the Gaussian
density and the calibrated slice measure cancels the rescaling factors from
Theorem~\ref{thm:m13-rescaling}; this is the reduced-volume transport identity.

For small-time stability, choose a gauge neighborhood $U$ of the base and
a smaller neighborhood whose spatial coordinates lie in a compact convex
box.  Smooth ODE dependence gives, uniformly for initial vectors in an open
neighborhood of $B\cup\{0\}$, candidate branches of action at most
$A\sqrt\tau$ and bounded coordinate speed in square-root time.  The
whole-slab curvature bound bounds the scalar term of every competitor,
including one outside $U$.  A competitor which leaves the smaller
neighborhood must traverse a fixed positive coordinate distance.  Its
kinetic cost is bounded below by a constant times $1/\sqrt\tau$, whereas
the negative scalar contribution is at most a constant times $\tau^{3/2}$.
For small $\tau$ it cannot beat the candidate's $A\sqrt\tau$ action.
This proves confinement of all relevant competitors rather than assuming it.

In the convex box let $m>0$ be a common metric lower bound, $K$ a coefficient
bound and $M$ the candidate speed bound.  The finite-energy comparison of a
candidate $p$ and a competitor $q$ with the same endpoints has a positive
leading term controlled by
\[
 \frac m4-\beta\tau>0,\qquad
 \beta=\frac{KM^2}{2}+K+\frac{(KM)^2}{m}.
\]
The error terms follow from the coefficient bounds and Cauchy--Schwarz;
the endpoint-zero difference satisfies
$\int|q-p|^2\leq\tau\int|(q-p)'|^2$ in the square-root interval.
Thus the comparison is nonnegative, and equality forces identical coordinate
curves.  Choose the common time interval smaller than the ODE, confinement
and positive-coefficient intervals.  This proves unique global minimality
for the entire open source neighborhood.  Uniform noncriticality on $B$
and compact capture then give the claimed stable sets.  Small-time coverage
is therefore a proved output of the generalized theory.
Compact exhaustion of the initial tangent space and the integrable initial
Gaussian give the zero-time volume limit by dominated convergence.

For ordinary capture the product cylinder must carry exactly the ordinary
metric.  A path projection preserves clock, horizontal speed and scalar
term; its inverse lift preserves admissibility and endpoints.  The action
sets and their infima therefore agree on the captured region.  The inverse
endpoint branches agree there as well, and the slice identification preserves
calibrated measure on measurable captured subsets.  Applying the ordinary
existence, differential and volume theorems supplies the three conclusions.
These theorems are explicit inputs to this capture operation, even though
they have already been constructed for ordinary complete bounded-curvature
flows.  The assertion makes no attainment claim for an arbitrary finite-action
generalized endpoint.
\end{proof}

\section{Noncollapsing}

\begin{definition}[Parabolic noncollapse]
\label{def:parabolic-noncollapse}
\lean{PoincareMT.M15GeneralizedNoncollapseAt}
\uses{def:generalized-spacetime}
For $\kappa,r_0>0$, a point $(x,t)$ is $\kappa$-noncollapsed through scale
$r_0$ if for every $0<r\leq r_0$ and every compatible cylinder based on
the entire terminal ball $B_{g(t)}(x,r)$, of time interval $[t-r^2,t]$,
whose full curvature is at most $r^{-2}$ on its whole image, including the
time endpoints, one has
\[
 \mu_{g(t)}(B_{g(t)}(x,r))\geq\kappa r^n
\]
The generalized predicate uses the actual slice measure
and actual carrier ball; a transported chart or a different carrier is not an
implicit replacement.
\end{definition}

\begin{theorem}[Generalized and compact noncollapsing]
\label{thm:m15-noncollapse}
\uses{def:parabolic-noncollapse,thm:m14-generalized}
\lean{PoincareMT.noncollapsingGeneralizedAndCompact}
For each $n>0$ and each $\bar\tau,l_0,V>0$ there is a constant
\[
 \kappa=\kappa(n,\bar\tau,l_0,V)>0
\]
with the following property for every generalized flow carrying the preceding
calculus.  Let $x\in X_T$, $r>0$, and let a compatible cylinder identify
its terminal source with exactly $B_{g(T)}(x,r)$ over $[T-r^2,T]$.
Assume full curvature at most $r^{-2}$ everywhere on its image and compact
closure of this terminal ball.  Suppose
\[
 r^2\leq\tau_0\leq\bar\tau,\qquad T-\tau_0\in I,
\]
and let $E$ be one exponential family based at $(x,T)$.  Let $W$ be an open
subset of its stable minimizing initial vectors at time $\tau_0$ such that
its normalized branch action is at most $l_0$ on $W$, and the image of $W$
in the slice $X_{T-\tau_0}$ has measure at least $V$.  Then
\[
 \mu_{g(T)}(B_{g(T)}(x,r))\geq\kappa r^n.
\]
In particular, if a region supplies such a configuration with the same
$(\bar\tau,l_0,V)$ for every tested point, radius and cylinder through
scale $r_0$, the whole region is $\kappa$-noncollapsed through that scale.

Separately, for every $\omega,T_0>0$ there is $\kappa_c(\omega,T_0)>0$
such that the following holds on every compact three-manifold, not necessarily
connected.  Let $g(t)$ be an ordinary Ricci flow on $[0,T]$, with
$0<T\leq T_0$, $|\mathrm{Rm}_{g(0)}|\leq1$, and
$\mu_{g(0)}(B_{g(0)}(q,1))\geq\omega$ for every $q$.
For $t_0\in[0,T]$ and $0<r$ with $r^2\leq t_0$, assume
\[
 |\mathrm{Rm}_{g(s)}(q)|\leq r^{-2}
 \quad\text{for all }s\in[t_0-r^2,t_0],\quad
                         q\in B_{g(t_0)}(p,r).
\]
Then $\mu_{g(t_0)}(B_{g(t_0)}(p,r))\geq\kappa_c r^3$.
The ordinary-flow results above supply the capture inputs for this assertion.
\source{MT2007}\dcref{Theorem 8.1 and Proposition 8.2, pp. 169--171}
\dcref{Lemma 8.3, Corollary 8.4, Claims 8.5--8.8 and Lemma 8.9, pp. 171--176}
\dcref{Theorem 8.10, pp. 176--177}\cite[Theorem 8.1, Proposition 8.2, Lemma 8.3, Corollary 8.4, Claims 8.5--8.8, Lemma 8.9, and Theorem 8.10]{MT2007}
\end{theorem}
\begin{proof}
First fix the constants, before any flow or cylinder is selected.  The
terminal assumptions and monotonicity on the same stable source $W$ give,
for every $0<\tau\leq\tau_0$,
\begin{equation}\label{eq:noncollapse-lower}
 V_{\mathrm{red}}(W,\tau)\geq
 \tau_0^{-n/2}e^{-l_0}V\geq
 c:=\bar\tau^{-n/2}e^{-l_0}V>0.
\end{equation}
Indeed the minimizing branch action realizes the infimum on the stable
image; its density at $\tau_0$ is at least
$\tau_0^{-n/2}e^{-l_0}$.  Calibrated change of variables preserves that
image volume.  The strict-prefix stable sets use the same $W$, so
monotonicity applies all the way through $\tau_0$.

We next prove the upper estimate used in the contradiction.  For a
dimension-only $\varepsilon_0\leq1/2$, suppose
$0<\varepsilon\leq\varepsilon_0$ and
$\mu_{g(T)}(B(x,r))\leq\varepsilon^n r^n$.
Evaluate at $\tau=\varepsilon r^2$ and split the measurable source into
\[
 W_1=W\cap\{|Z|^2\leq(64\varepsilon)^{-1}\},\qquad
 W_2=W\setminus W_1.
\]
For $W_1$, the cylinder estimates confine its small-initial-vector curves
inside the actual cylinder.  Curvature bounds give
$R\geq-n^2r^{-2}$ along them and compare the measure of their endpoint
image at time $T-\varepsilon r^2$ with
$e^{n^2\varepsilon}\mu_{g(T)}(B(x,r))$.  Kinetic energy is nonnegative,
so the scalar bound gives $l\geq-n^2\varepsilon/3$.  Integration over
that endpoint image gives
\[
 V_{\mathrm{red}}(W_1,\varepsilon r^2)
 \leq (\varepsilon r^2)^{-n/2}
       e^{n^2\varepsilon/3}e^{n^2\varepsilon}\varepsilon^nr^n
 \leq2\varepsilon^{n/2},
\]
after shrinking $\varepsilon_0$.  The radius factors have cancelled exactly.
For completeness, the confinement has a strict interior margin.  The
initial square-root speed is $2|Z|\leq1/(4\sqrt\varepsilon)$.
For a lifted prefix staying in the terminal half-ball and
$0\leq s^2\leq r^2/2$, the interior scalar-gradient estimate has the form
$|\nabla R|\leq A/r^3$, where $A$ depends only on dimension.
To apply Theorem~\ref{thm:rf-shi}, rescale the same cylinder by $r^{-2}$
to $[0,1]$ and unit curvature.  At a point of its terminal half-ball,
metric comparison puts the initial-metric ball of radius $e^{-n}/8$
inside a compact terminal interior buffer.  Indeed an initial path of that
length has terminal length at most $1/8$, leaving a positive margin inside
the terminal unit ball.  Compact closure of the tested ball then gives
the required compact initial ball.  The local derivative estimate applies
at its center for times in $[1/2,1]$ and gives a dimension-only bound for
$|\nabla R|$.  Scaling back gives precisely $A/r^3$.
Pairing the square-root Euler equation with its velocity,
using $|\operatorname{Ric}(v,v)|\leq nr^{-2}|v|^2$, and integrating the
resulting scalar differential inequality gives, after metric comparison to
the terminal slice,
\[
 |\alpha'(s)|_{g(T)}\leq
 e^{2ns^2/r^2}\left(|\alpha'(0)|+\frac{2A}{3r^3}s^3\right).
\]
For $s\leq\sqrt\varepsilon r$ choose the dimension-only threshold so
that this is at most $5/(12\sqrt\varepsilon)$.
Integration gives displacement at most $5r/12$, strictly less than $r/2$.
The inverse cylinder map gives a lift initially.
If its maximal confined prefix stopped before the required time, its image
would lie in the compact closed $5r/12$ buffer, strictly inside the source
half-ball.  Finitely many inverse chart neighborhoods then extend the lift
beyond that endpoint, a contradiction.  Thus the entire prefix remains
in the cylinder.  This establishes the scalar and endpoint-volume estimates
on the actual family, rather than on a separately chosen comparison curve.

For $W_2$ use the initial Gaussian bound instead.  If $a=(64\varepsilon)^{-1}$,
then on $W_2$
$e^{-|Z|^2}\leq e^{-a/2}e^{-|Z|^2/2}$, whence
\[
 V_{\mathrm{red}}(W_2,\varepsilon r^2)
 \leq 2^n(2\pi)^{n/2}e^{-1/(128\varepsilon)}
 \leq\varepsilon^{n/2}.
\]
The last inequality holds below another dimension-only threshold because
exponential decay dominates every positive power.  Disjoint stable images
have additive reduced volume, so the total is at most
$3\varepsilon^{n/2}$.  Choose $\varepsilon$ below both thresholds with
$3\varepsilon^{n/2}<c$ and set $\kappa=\varepsilon^n$.  A violating
ball would satisfy the upper estimate, whereas
$\varepsilon r^2\leq r^2\leq\tau_0$ allows
\eqref{eq:noncollapse-lower}; this is the contradiction.  Notice that the
positive number in the contradiction is $c$, not the unweighted volume $V$.
The pointwise-region conclusion follows by applying the same constant to
each supplied configuration.

For the compact theorem let $\delta=(32\cdot3^6)^{-1}$.
Up to time $\delta$, the initial full-curvature bound and curvature evolution
give the short-time doubling bound.  Integrating metric and volume evolution
and applying initial small-ball comparison gives
$\mu_{g(t_0)}(B(p,r))\geq k_e(\omega)r^3$ for $t_0<\delta$;
here $r^2\leq t_0$ implies $r\leq1$.
For $t_0\geq\delta$, set $b=t_0-\delta/4$.
The reduced-length minimum bound, followed by a short comparison path in
the initial controlled region, gives a point $q_0$ whose initial unit ball
has a uniform reduced-length bound at backward time $b$.
More explicitly choose $q_0$ at backward time $c=t_0-\delta/2$ with
$l(q_0,c)\leq3/2$, and join it to a point of its initial unit ball over
$[c,b]$.  On the corresponding physical interval $[\delta/4,\delta/2]$,
the curvature bound gives $R\leq18$ and norm comparison factor
$C=e^{6\delta}$.  A short initial-metric path of speed at most
$C/(b-c)$ has additional action at most
$\sqrt b(18+(C/(b-c))^2)(b-c)$.  Division by $2\sqrt b$ gives the
uniform choice
\[
 l_0=\frac32+\frac{\delta}{8}
              \left(18+\left(\frac{C}{\delta/4}\right)^2\right).
\]
Smooth concatenations approximate the joined path with arbitrarily small
extra action, which suffices for the infimum inequality.
Initial volume comparison gives a positive lower bound $V(\omega)$ for
the measure of this ball in the early slice $g(\delta/4)$.
The ordinary capture map has a full-measure stable image, so its inverse
image of this ball supplies the open set $W$ and preserves this lower bound.
This constructs a generalized configuration, rather than postulating one.

Use the tested cylinder with radius $r/2$.  It retains the curvature bound
and satisfies $(r/2)^2\leq b$.  Apply the generalized estimate with
$\bar\tau=T_0$, the uniform reduced-length bound and $V(\omega)$.
Transport the slice measure back through the same ordinary-product
identification and include the half-radius ball in the original ball:
the resulting constant is $k_l=\kappa/8$.
Finally choose $\kappa_c=\min(k_e,k_l)>0$ before the manifold and flow.
For a disconnected compact manifold restrict to the component of $p$;
the tested balls and their measures agree with those in the ambient manifold.
\end{proof}

\section{Contracts and limitations}

The outputs of this chapter are the monotonicity and equality-rigidity theorem
(Theorem~\ref{thm:m10-reduced-volume}), the generalized carrier/gauge and
rescaling identities (Theorems~\ref{thm:m11-m12-gauge} and
\ref{thm:m13-rescaling}), generalized reduced-volume transport
(Theorem~\ref{thm:m14-generalized}), and the two noncollapse interfaces
(Theorem~\ref{thm:m15-noncollapse}).  The ancient-models chapter may consume
the generalized cylinder and its noncollapse conclusion, but must preserve
the carrier and scale sequence.  The Ricci-flow chapter supplies the curvature
theory and the ordinary slab hypotheses.  The singular-limits chapter consumes the
generalized carrier, exact rescaling, and pointwise noncollapse; it does not
receive completeness for free.

The generalized estimate is conditional on its explicitly stated stable-set
configuration.  A later application must construct that configuration on
its own flow and with constants uniform over the tested region.  The compact
case above supplies one such construction; it does not supply it for every
generalized limit.  The curvature inputs used here are the scalar and metric
evolution identities and Theorem~\ref{thm:rf-shi}; the compact case also uses
the short-time curvature bound and initial relative ball-volume comparison.
Their hypotheses must be preserved when these results are applied on a gauge
cylinder or a component of an ordinary flow.
